%-----------------------------------------------------------------------------------------------------------------------------------------------%
%	The MIT License (MIT)
%
%	Copyright (c) 2015 Jan Küster
%
%	Permission is hereby granted, free of charge, to any person obtaining a copy
%	of this software and associated documentation files (the "Software"), to deal
%	in the Software without restriction, including without limitation the rights
%	to use, copy, modify, merge, publish, distribute, sublicense, and/or sell
%	copies of the Software, and to permit persons to whom the Software is
%	furnished to do so, subject to the following conditions:
%
%	THE SOFTWARE IS PROVIDED "AS IS", WITHOUT WARRANTY OF ANY KIND, EXPRESS OR
%	IMPLIED, INCLUDING BUT NOT LIMITED TO THE WARRANTIES OF MERCHANTABILITY,
%	FITNESS FOR A PARTICULAR PURPOSE AND NONINFRINGEMENT. IN NO EVENT SHALL THE
%	AUTHORS OR COPYRIGHT HOLDERS BE LIABLE FOR ANY CLAIM, DAMAGES OR OTHER
%	LIABILITY, WHETHER IN AN ACTION OF CONTRACT, TORT OR OTHERWISE, ARISING FROM,
%	OUT OF OR IN CONNECTION WITH THE SOFTWARE OR THE USE OR OTHER DEALINGS IN
%	THE SOFTWARE.
%
%
%-----------------------------------------------------------------------------------------------------------------------------------------------%


%============================================================================%
%
%	DOCUMENT DEFINITION
%
%============================================================================%

%we use article class because we want to fully customize the page and dont use a cv template
\documentclass[10pt,A4]{article}


%----------------------------------------------------------------------------------------
%	ENCODING
%----------------------------------------------------------------------------------------

%we use utf8 since we want to build from any machine
\usepackage[utf8]{inputenc}

%----------------------------------------------------------------------------------------
%	LOGIC
%----------------------------------------------------------------------------------------

% provides \isempty test
\usepackage{xifthen}
\usepackage{enumitem}

%----------------------------------------------------------------------------------------
%	FONT
%----------------------------------------------------------------------------------------

% some tex-live fonts - choose your own

%\usepackage[defaultsans]{droidsans}
%\usepackage[default]{comfortaa}
%\usepackage{cmbright}
\usepackage[default]{raleway}
%\usepackage{fetamont}
%\usepackage[default]{gillius}
%\usepackage[light,math]{iwona}
%\usepackage[thin]{roboto}

% set font default
\renewcommand*\familydefault{\sfdefault}
\usepackage[T1]{fontenc}

% more font size definitions
\usepackage{moresize}


%----------------------------------------------------------------------------------------
%	PAGE LAYOUT  DEFINITIONS
%----------------------------------------------------------------------------------------

%debug page outer frames
%\usepackage{showframe}


%define page styles using geometry
\usepackage[a4paper]{geometry}

% for example, change the margins to 2 inches all round
\geometry{top=2cm, bottom=1cm, left=1.5cm, right=1.5cm}

%use customized header
\usepackage{fancyhdr}
\pagestyle{fancy}

%less space between header and content
\setlength{\headheight}{-5pt}


%customize entries left, center and right
%\lhead{}
%\chead{}
%\rhead{}


%indentation is zero
\setlength{\parindent}{0mm}

%----------------------------------------------------------------------------------------
%	TABLE /ARRAY DEFINITIONS
%----------------------------------------------------------------------------------------

%for layouting tables
\usepackage{multicol}
\usepackage{multirow}

%extended aligning of tabular cells
\usepackage{array}

\newcolumntype{x}[1]{%
>{\raggedleft\hspace{0pt}}p{#1}}%


%----------------------------------------------------------------------------------------
%	GRAPHICS DEFINITIONS
%----------------------------------------------------------------------------------------

%for header image
\usepackage{graphicx}

%for floating figures
\usepackage{wrapfig}
\usepackage{float}
%\floatstyle{boxed}
%\restylefloat{figure}

%for drawing graphics
\usepackage{tikz}
\usetikzlibrary{shapes, backgrounds,mindmap, trees}


%----------------------------------------------------------------------------------------
%	Color DEFINITIONS
%----------------------------------------------------------------------------------------

\usepackage{color}

%accent color
%\definecolor{sectcol}{RGB}{0,150,255}
\definecolor{sectcol}{RGB}{0,0,0}

%dark background color
\definecolor{bgcol}{RGB}{110,110,110}

%light background / accent color
\definecolor{softcol}{RGB}{205,205,205}

\usepackage{hyperref}
\usepackage{xcolor}
\hypersetup{
    colorlinks,
    linkcolor={red!50!black},
    citecolor={blue!50!black},
    urlcolor=[RGB]{0 150 255}
}

%============================================================================%
%
%
%	DEFINITIONS
%
%
%============================================================================%

%----------------------------------------------------------------------------------------
% 	HEADER
%----------------------------------------------------------------------------------------

% remove top header line
\renewcommand{\headrulewidth}{0pt}

%remove botttom header line
\renewcommand{\footrulewidth}{0pt}

%remove section num
\renewcommand{\thesection}{}

%----------------------------------------------------------------------------------------
% 	ARROW GRAPHICS in Tikz
%----------------------------------------------------------------------------------------

% a six pointed arrow poiting to the left
\newcommand{\tzlarrow}{(0,0) -- (0.2,0) -- (0.3,0.2) -- (0.2,0.4) -- (0,0.4) -- (0.1,0.2) -- cycle;}

% include the left arrow into a tikz picture
% param1: fill color
%
\newcommand{\larrow}[1]
{\begin{tikzpicture}[scale=0.58]
	 \filldraw[fill=#1!100,draw=#1!100!black]  \tzlarrow
 \end{tikzpicture}
}

% a six pointed arrow poiting to the right
\newcommand{\tzrarrow}{ (0,0.2) -- (0.1,0) -- (0.3,0) -- (0.2,0.2) -- (0.3,0.4) -- (0.1,0.4) -- cycle;}

% include the right arrow into a tikz picture
% param1: fill color
%
\newcommand{\rarrow}
{\begin{tikzpicture}[scale=0.7]
	\filldraw[fill=sectcol!100,draw=sectcol!100!black] \tzrarrow
 \end{tikzpicture}
}



%----------------------------------------------------------------------------------------
%	custom sections
%----------------------------------------------------------------------------------------

% create a coloured box with arrow and title as cv section headline
% param 1: section title
%
\newcommand{\cvsection}[1]
{
	\begin{flushleft}
  \LARGE\textcolor{sectcol}{\textbf{  #1}}
	\end{flushleft}
  \vspace{-15pt}
  \textcolor{softcol}{\hrule}
  \vspace{8pt}
}

%create a coloured arrow with title as cv meta section section
% param 1: meta section title
%
\newcommand{\metasection}[3]
{
%\begin{tabular*}{1\textwidth}{r r}
%\footnotesize{#1} \hspace*{\fill} \footnotesize{#2}\hspace*{\fill} \footnotesize{#3}\\[1pt]
\small{#1} \hspace*{\fill} \small{#2}\hspace*{\fill} \small{#3}\\[1pt]
%\end{tabular*}
}

%----------------------------------------------------------------------------------------
%	 CV EVENT
%----------------------------------------------------------------------------------------

% creates a stretched box as cv entry headline followed by some paragraphs about
% the work you did
% param 1:	event time i.e. 2014 or 2011-2014 etc.
% param 2:	event name (what did you do?)
% param 3:	institution (where did you work / study)
% param 4:	list of paragraphs
%
\newcommand{\cvevent}[4]
{

\begin{tabular*}{1\textwidth}{p{13.6cm}  x{3.9cm}}
\textbf{#2} - \textcolor{bgcol}{#3} &   \vspace{2.5pt}\textcolor{sectcol}{#1}
\end{tabular*}

%\textcolor{softcol}{\hrule}
\vspace{-4pt}
\begin{itemize}
 \setlength\itemsep{-0.05cm}
	\foreach \desc in {#4}{
    %$\cdot$ \desc\\[3pt]
    \item \desc
  }
\end{itemize}
\vspace{3pt}
}

\newcommand{\cveventsum}[5]
{

\begin{tabular*}{1\textwidth}{p{13.6cm}  x{3.9cm}}
\textbf{#2} - \textcolor{bgcol}{#3} &   \vspace{2.5pt}\textcolor{sectcol}{#1}
\end{tabular*}

\vspace{-8pt}
\begin{description}[leftmargin=*,labelindent=1pt]
 \item #4
\end{description}
%\textcolor{softcol}{\hrule}
\vspace{-10pt}
\begin{itemize}
 \setlength\itemsep{-0.05cm}
	\foreach \desc in {#5}{
    %$\cdot$ \desc\\[3pt]
    \item \desc
  }
\end{itemize}
\vspace{3pt}
}

\newcommand{\cveventdot}[4]
{

\begin{tabular*}{1\textwidth}{p{13.6cm}  x{3.9cm}}
\textbf{#2} - \textcolor{bgcol}{#3} &   \vspace{2.5pt}\textcolor{sectcol}{#1}
\end{tabular*}

\begin{description}[labelindent=0cm]
	\foreach \desc in {#4}{
   \vspace{-6pt}
   \item \desc
  }
\end{description}

\vspace{3pt}
}

\newcommand{\cveventtwo}[7]
{

%\begin{tabular}{2\textwidth}{p{13.6cm} x{3.9cm}}
\begin{tabular*}{2\textwidth}{p{13.6cm} c}
\textbf{#2} - \textcolor{bgcol}{#3} &   \vspace{2.5pt}\textcolor{sectcol}{#1} \\
\textbf{#5} - \textcolor{bgcol}{#6}  &   \vspace{2.5pt}\textcolor{sectcol}{#4}
\end{tabular*}

%\textcolor{softcol}{\hrule}
\vspace{-4pt}
\begin{itemize}
 \setlength\itemsep{-0.05cm}
	\foreach \desc in {#7} {
		%$\cdot$ \desc\\[3pt]
    \item \desc
	}
\end{itemize}

\vspace{-2pt}
\vspace{3pt}
}

\newcommand{\cveventtwosum}[8]
{

%\begin{tabular}{2\textwidth}{p{13.6cm} x{3.9cm}}
\begin{tabular*}{2\textwidth}{p{13.6cm} c}
\textbf{#2} - \textcolor{bgcol}{#3} &   \vspace{2.5pt}\textcolor{sectcol}{#1} \\
\textbf{#5} - \textcolor{bgcol}{#6}  &   \vspace{2.5pt}\textcolor{sectcol}{#4}
\end{tabular*}

\vspace{-8pt}
\begin{description}[leftmargin=*,labelindent=1pt]
 \item #7
\end{description}
\vspace{-6pt}

%\textcolor{softcol}{\hrule}
\vspace{-4pt}
\begin{itemize}
 \setlength\itemsep{-0.05cm}
	\foreach \desc in {#8} {
		%$\cdot$ \desc\\[3pt]
    \item \desc
	}
\end{itemize}

\vspace{-2pt}
\vspace{3pt}
}

\newcommand{\cveventtwodot}[7]
{

%\begin{tabular}{2\textwidth}{p{13.6cm} x{3.9cm}}
\begin{tabular*}{2\textwidth}{p{16.7cm} c}
\textbf{#2} - \textcolor{bgcol}{#3} &   \vspace{2.5pt}\textcolor{sectcol}{#1} \\
\textbf{#5} - \textcolor{bgcol}{#6}  &   \vspace{2.5pt}\textcolor{sectcol}{#4}
\end{tabular*}

%\textcolor{softcol}{\hrule}
\vspace{4pt}

	\foreach \desc in {#7} {
		$\cdot$ \desc\\[3pt]
	}

\vspace{-2pt}
\vspace{3pt}
}





% creates a stretched box as
\newcommand{\cveventmeta}[2]
{
	\mbox{\mystrut \hspace{87pt}\textit{#1}}\\
	#2
}

%----------------------------------------------------------------------------------------
% CUSTOM STRUT FOR EMPTY BOXES
%----------------------------------------- -----------------------------------------------
\newcommand{\mystrut}{\rule[-.3\baselineskip]{0pt}{\baselineskip}}

%----------------------------------------------------------------------------------------
% CUSTOM LOREM IPSUM
%----------------------------------------------------------------------------------------
\newcommand{\lorem}
{Lorem ipsum dolor sit amet, consectetur adipiscing elit. Donec a diam lectus.}

%Set header and footerd
%\fancyhf{} % sets both header and footer to nothing
\renewcommand{\headrulewidth}{0pt}
% your new footer definitions here
\lfoot{weiweic}
\cfoot{résumé -- \today}
\rfoot{\thepage}

%============================================================================%
%
%
%
%	DOCUMENT CONTENT
%
%
%
%============================================================================%
\begin{document}


%---------------------------------------------------------------------------------------
%	TITLE HEADLINE
%----------------------------------------------------------------------------------------
\vspace{-8pt}
\begin{center}
	%\HUGE \textsc{Weiwei Chen}
	\HUGE Weiwei Chen
\end{center}

\vspace{6pt}


%---------------------------------------------------------------------------------------
%	META SECTION
%----------------------------------------------------------------------------------------
\metasection{\textbf{Homepage:} \href{https://weiweichen.github.io}{weiweichen.github.io}}{\textbf{Email:} \href{mailto:weiwei.chen.uci@gmail.com}{weiwei.chen.uci@gmail.com}}
{\textbf{} }
%{\textbf{Phone Number:} (949) 887-6878}
\vspace{-2pt}
\vspace{6pt}

\normalsize

%---------------------------------------------------------------------------------------
% Research Interests
%----------------------------------------------------------------------------------------
\vspace{-6pt}
\cvsection{Research Interests}

\vspace{-4pt}
\begin{itemize}[leftmargin=*]
 \setlength\itemsep{-0.5cm}
  \item Compilers and Frameworks for Machine Learning, High-Performance Computing, and Big Data \\
  \item Heterogeneous Parallel Programming
\end{itemize}
\vspace{8pt}

%============================================================================%
%
%	CV SECTIONS AND EVENTS (MAIN CONTENT)
%
%============================================================================%

%---------------------------------------------------------------------------------------
%	EDUCATION SECTION
%--------------------------------------------------------------------------------------
\vspace{-12pt}
\cvsection{Education}

\cveventdot{2013}{Ph.D. Electrical and Computer Engineering}{University of California, Irvine}{
   {Dissertation: Out-of-Order Parallel Discrete Event Simulation for Electronic System Level Design},
   {Outstanding Dissertation Award, European Design and Automation Association}
}
\vspace{-8pt}
\cveventtwodot{2007}{M.S. Computer Engineering}{Shanghai Jiao Tong University}
{2004}{B.S. Computer Science and Engineering}{Shanghai Jiao Tong University} {
}
\vspace{-30pt}

%---------------------------------------------------------------------------------------
%	EXPERIENCE
%----------------------------------------------------------------------------------------
\cvsection{Experience}
\cveventtwosum{March 2025 - Present}{Senior  Staff Engineer}{Modular}
  {May 2023 - Feb 2025}{Staff Engineer}{Modular}
  {Building next-generation AI infrastructure and develop platform}
  {
  {Compiler for the \href{https://www.modular.com/mojo}{Mojo} programming language.}
}

\cveventtwosum{May 2021 - May 2023}{Senior Principal Engineer}{SambaNova Systems}
        {Oct 2018 - May 2021}{Principal Engineer}{SambaNova Systems}
        {{Building Compilers for Coarse-Grain Reconfigurable Dataflow Units (RDUs)}}
        {
  {\textbf{Arc} - An MLIR-based DSL and Compiler for stitching high performance tensor operation templates on RDUs,
    including Layout and Control Analysis, Buffer and Anti-hang Peephole Optimizations, Template Selection and Lowering, etc.
    One of the first two engineers brought up this layer.},
  {\textbf{RAIL} - An MLIR-based DSL and Compiler for programming the compute and memory components on RDUs.
  Founding member of the project and main contributor to the compute component, including syntax design, register allocation,
  context splitting, latency analysis, etc.},
  {\textbf{Nova} - PyTorch JIT for RDU, including adding dynamic dispatch key for RDU, RDU compiler integration using lazy tensor, hot/cold caching mechanisms,
  performance annotation for auto-differentiated graphs, heterogeneous backend enabling for CPU+GPU+RDU, etc.
  One of the first two engineers bringing up this layer.},
  {\textbf{Assembler} for multi-generation RDUs, founding member of the project to design and implement an in-house data structure
  for the memory component and assembler in general.},
  {Various compiler infrastructure tasks, including LLVM/MLIR version upgrading; architecting and migrating one compiler layer from Python to MLIR in C++;
  migrating RDU backend from Python to C++; cmake structuring and maintenance; ML model compilation debugging support, etc},
  {Various compiler PoC projects, including a \href{https://spatial-lang.org}{Spatial}-like RDU compiler in Scala; JAX to RDU lowering;
  new DSL for RDU programming at higher abstraction level than RAIL, etc.}
}

\cveventsum{Feb 2016 - Oct 2018}{Founding Software Engineer}{BigStream Solutions}
  {First Engineer hired to lead software research and development efforts on compilers,
   native C++ acceleration libraries, tool-chains for FPGA acceleration, and big data system architecture.}
  {
  {A dataflow compiler in Scala for Spark SQL query to intermediate representations with three backend for
   native CPU, FPGA, and RISCV code generation, and a planner for optimized SW/HW accelerator partitioning.},
  {Native C++ runtime and compiler support for accelerating user-defined functions(UDF) in Scala.},
  {A C++ native acceleration library with templatized SQL operations, cluster data source support,
   (HDFS, amazon S3, Microsoft WASB), various input format support (json, avro, parquet, csv)
   and TensorFlow native integration for data pipelines.},
  {A Clang-LLVM based high-level synthesis compiler for timing scheduling and code generation for FPGA.},
  {Various engineering support, including on-premise Hadoop Cluster setup, network performance measurement, etc.}
}

\pagebreak

\cveventsum{Oct 2013 - Feb 2016}{Senior Engineer}{Qualcomm Research Silicon Valley}
 {Member of the Pervasive Concurrency team for Qualcomm Symphony System Manager SDK, a task-based parallel programming patterns and runtime for heterogeneous multi-core platforms.}
{{Heterogeneous parallel pipeline programming pattern API and internal scheduling, task and dataflow API infrastructures,
parallelize Android native computational photography and enterprise compression applications using task-based parallel patterns,
power and performance evaluation for native parallel applications, etc.},
{Compiler frontend analysis and backend code generation for coarse-grain auto-parallelization base on polyhedral optimizations in LLVM.}
}

\cveventsum{June 2011 - Sept 2011}{Software Develop Engineer Intern}{Microsoft}
 {Windows Core Security and Identity Public Key Infrastructure Team}
{
 {A Windows store application for secure banking on Windows 8.}
}

%---------------------------------------------------------------------------------------
%	Patents SECTION
%--------------------------------------------------------------------------------------
\cvsection{Patents (selected)}
\vspace{-4pt}
\begin{itemize}
 \setlength\itemsep{-0.1cm}
 \item Iteration Synchronization Construct for Parallel Pipelines (granted), \textcolor{bgcol}{US 15/191,266}
 \item Systems And Methods For Accelerating Data Operations By Utilizing Native Memory Management (filed),\textcolor{bgcol}{US 62/775,533}
 \item Systems and Methods for Accelerating Data Operations By Utilizing Dataflow Subgraph Templates (filed), \textcolor{bgcol}{US 16/898,048}
 \item Compile Time Logic for Inserting a Buffer Between a Producer Operation Unit and a Consumer Operation Unit in a Dataflow Graph (granted), \textcolor{bgcol}{US 17/582,421}
 \item Compile Time Logic for Detecting Streaming Compatible and Broadcast Compatible Data Access Patterns (filed), \textcolor{bgcol}{US 17/031,679}
 \item Anti-Congestion Flow Control for Reconfigurable Processors (filed), \textcolor{bgcol}{US 16/890,841}
 \item Partitioning Dataflow Operations for a Reconfigurable Computing System (filed), \textcolor{bgcol}{US 63/317,476}
\end{itemize}

%---------------------------------------------------------------------------------------
%	Publications SECTION
%--------------------------------------------------------------------------------------
\cvsection{Publications (selected)}
\vspace{-4pt}
\begin{itemize}
 \setlength\itemsep{-0.1cm}
\item Weiwei Chen, Xu Han, Rainer Doemer, \textcolor{bgcol}{"May-Happen-in-Parallel Analysis based on Segment Graphs for Safe ESL Models"},
  in Proceedings of the \emph{Design, Automation and Test in Europe Conference (DATE)}, Dresden, Germany, March 2014 \textbf{(Best Paper Award)}
 \item Weiwei Chen, Xu Han, Che-Wei Chang, Guantao Liu, Rainer Doemer, \textcolor{bgcol}{"Out-of-Order Parallel Discrete Event Simulation for Transaction Level Models"},
  \emph{IEEE Transactions on Computer-Aided Design of Integrated Circuits and Systems (TCAD)}, vol.33, no.12, pp.1859-1872, Dec 2014
 \item 3 Journal papers, 1 Book, 2 Book Chapters, 14 Conference papers, 4 Technical Reports, 4 Poster Presentations,
  see full list from \href{http://www.cecs.uci.edu/~weiweic/pubs\_by\_type.html}{here}.
\end{itemize}

%---------------------------------------------------------------------------------------
% Awards and Honors
%--------------------------------------------------------------------------------------
\cvsection{Awards and Honors (selected)}
\vspace{-4pt}
\begin{itemize}
 \setlength\itemsep{-0.1cm}
 \item 7 Qualstar Awards, Qualcomm Inc. 2014 - 2016
 \item Outstanding Dissertation Award, European Design and Automation Association (EDAA) 2014
 \item Best Paper Award, Design, Automation and Test Conference in Europe (DATE) 2014
 \item Pedagogical Fellowship, UC Irvine 2012-13
 \item Henry Samueli Endowed Fellowship, UC Irvine 2007
 \item National Scholarship for Academic Excellence, China 2006
 \item Exceptional Undergraduate Student Awards, SJTU
 \item People's Scholarship for Academic Excellence, SJTU 2000-2004
\end{itemize}

%---------------------------------------------------------------------------------------
% Skills and Hobbies
%--------------------------------------------------------------------------------------
\cvsection{Skills and Hobbies}
\vspace{-4pt}
\begin{itemize}
 \setlength\itemsep{-0.1cm}
 \item C++ (native), Scala (product), Python (project), Rust (Advent of Code), Haskell (Advent of Code),  Git, PyTorch, Vim, GDB, Valgrind, cmake, Latex ...
 \item Board member of San Francisco Bay Queer Contra Dance
 \item Fiddling, Reading, Hiking, Folk Dancing, Gardening, Karate, Sailing, Board Games
\end{itemize}

\cvsection{Working Status}
\vspace{-4pt}
\begin{itemize}
 \item US Citizen
\end{itemize}


%--------------------------------------------------------------------------------------------------
%	ARTIFICIAL FOOTER (fancy footer cannot exceed linewidth)
%--------------------------------------------------------------------------------------------------

%\null
%\vspace*{\fill}
%\hspace{-0.25\linewidth}\colorbox{white}{\makebox[1.5\linewidth][c]{\mystrut  \textnormal{\textcolor{sectcol}{www.jankuester.com} $\cdot$ \textcolor{sectcol}{github.com/jankapunkt}}}}




%============================================================================%
%
%
%
%	DOCUMENT END
%
%
%
%============================================================================%
\end{document}
